\chapter{Solutions to the checks}\label{app:solutions}

\section{Stable pointed curves}

\MCsolution{ex:stability-test}
The normalization has two points above the self-node and one
marked point, so $r=3$. For a rational component,
$2h-2+r=-2+3=1>0$. Counting the node only once would give
$r=2$, which would incorrectly declare the component unstable.
The local dualizing contribution comes from both branches.

\MCsolution{ex:total-vs-fiber}
The base map is $i:\Spec\CC\to\mathbb A^1$ at $0$ and the
total-space map is $C_0\to C_0\times\mathbb A^1$ at $0$.
The square is a fiber-product square, so this is a moduli arrow.
If it had an inverse in the total category, applying $p$ would
give an inverse to $i$, which does not exist. A fiber category
over $S$ contains only arrows over $\id_S$. The arrow just
constructed lies over $i$, not over an identity, so it is not
a counterexample to the groupoid assertion.

\MCsolution{ex:cartesian-uniqueness}
Apply the universal property of $\xi_T\to\xi$ to
$\xi'_T\to\xi$ with base factorization $\id_T$ to obtain
$b:\xi'_T\to\xi_T$. Interchanging the roles gives
$c:\xi_T\to\xi'_T$. Both $bc$ and the identity are
factorizations of the same arrow through $\xi_T\to\xi$,
so uniqueness gives $bc=\id$; similarly $cb=\id$.
The commuting condition determines this isomorphism. Without
that condition, the identity and a hyperelliptic involution
can be two different isomorphisms between the same genus-two
objects.

\MCsolution{ex:iso-sheaf}
Morphism descent produces a map, whose invertibility still
needs verification. Glue the compatible local inverse maps
and check the two composite maps against the identities
after the covers of the two total spaces. Uniqueness of
morphism gluing then makes them identities globally.
For each label $r$, compare the maps $a\sigma_r$ and
$\tau_r$ from the base to the target total space. Their
local equality gives global equality. Invertibility alone
would not say that labels are preserved.

\MCsolution{desc:exercise-polarization}
The canonical base-change isomorphisms for $\omega$ and
preservation of every marking divisor give
$a_{ij}^*L_j\simeq L_i$. Denote this isomorphism by
$\lambda_{ij}$. Functoriality carries the equation
$a_{ik}=a_{jk}a_{ij}$ to
\[
             \lambda_{ik}=\lambda_{ij}\circ a_{ij}^*\lambda_{jk}.
\]
Thus the bundle transitions have the required cocycle.
Unrelated ample bundles need not even pull back to isomorphic
bundles under the specified $a_{ij}$; and arbitrary chosen
isomorphisms, when available, need not satisfy a cocycle.
The canonical bundle supplies both existence and coherence
without adding a new choice to the moduli object.

\MCsolution{ex:twisted-datum}
An overlap point is a pair of square roots $u_1,u_2$ of the
same nonzero $t$. Either $u_2=u_1$ or $u_2=-u_1$.
Only the first comparison is on the diagonal. Two sign
changes compose to the identity since $\iota^2=\id$;
in coordinates $(u_2/u_3)(u_1/u_2)=u_1/u_3=1$.
From $y_i=z/u_i$ one gets
$y_j=(u_i/u_j)y_i$, exactly the prescribed transition map.
Checking only the diagonal of this overlap would lose all
the information about the involution.

\MCsolution{ex:curve-dm}
For a rational normalization component, stability gives $m_v\geq3$,
so $\deg T_{\mathbb P^1}(-B_v)=2-m_v<0$.
For an elliptic component, $m_v\geq1$ and the degree is $-m_v<0$.
For $h_v\geq2$, the degree $2-2h_v-m_v$ is already negative
even if $m_v=0$. An unmarked elliptic curve instead has trivial
tangent bundle and a nonzero translation vector field. A finite
group may have several elements: the hyperelliptic involution of
a smooth genus-two curve gives such an example. Triviality means
that only the identity remains, which is a stronger condition
than the one needed for a Deligne--Mumford stack.

\MCsolution{ex:curve-dimension}
There are two rational normalization components, each with three
special points, so both are stable. The genus is
$0+0+3-2+1=2$. Each component contributes
$3\cdot0-3+3=0$ parameters: a three-pointed rational curve has
no moduli. Each of the three nodes contributes one smoothing
parameter. The deformation space therefore has dimension $3$,
in agreement with $3g-3=3$.


\section{Stable maps}
\subsection*{Solutions for Chapter 2}

\MCsolution{maps:ex-stability}
A contracted rational component with no special points violates the requirement of three special points. Its canonical polarization has degree $-2$. A constant map from a smooth genus-two curve is stable because its only component has genus at least two; its polarization has degree $2g-2=2$. The pullback of $A$ under a constant map has degree zero in both cases.

\medskip
\MCsolution{maps:ex-arrow}
The identity of $\PP^1$ does not give this arrow: it would require $h\circ\operatorname{id}=f$, whereas the two displayed parametrizations are different maps. The coordinate swap is the required arrow. An arrow over a noninvertible base map cannot be invertible in the total category, since any inverse would project to an inverse of the base map. Its comparison with the base-changed target is nevertheless an isomorphism.

\medskip
\MCsolution{maps:ex-cartesian}
The total-space map must lie over $w$, preserve the marked sections, preserve the map to $X$, and induce an isomorphism with the base-changed target family. For the last target equation, $f_T=f\operatorname{pr}_C$ and $\operatorname{pr}_C\widetilde b=b$, so $f_T\widetilde b=fb=e$. The product description $E\simeq C\times_S U\simeq(C\times_S T)\times_T U$ proves the cartesian-square condition.

\medskip
\MCsolution{maps:ex-isom}
The compatible source isomorphisms glue to a unique isomorphism of the source schemes over $S$, and to a pointed-curve isomorphism if the markings were respected locally. Without $h_i a_i=f_i$ there is no reason for $ha=f$. The result would therefore be an arrow in a different moduli problem. For example, the identity of a source $\PP^1$ does not identify two different parametrizations of a line as stable maps unless the maps themselves agree.

\medskip
\MCsolution{maps:ex-descent}
Projectivity supplies the very ample $A$ and hence the relatively ample $\omega(\sum\sigma_i)\otimes f^*A^3$. The cocycle equation makes the curve and line-bundle identifications consistent on triple overlaps, so that the polarized descent theorem applies. Compatibility with $f_i$ identifies the $f_i^*A^3$ factors and later lets the maps to $X$ themselves glue. Underlying pointed curves need not be stable: an unmarked degree-one $\PP^1$ is already a counterexample. It is stable as a map because its nonconstant map contributes positivity and removes source automorphisms preserving the map.

\medskip
\MCsolution{maps:ex-tree}
The dual graph is a tree with exactly one vertex carrying positive degree. If it has at least two vertices, at least one leaf has degree zero. That contracted leaf has a single node and no markings, contradicting the three-special-point requirement for a contracted rational component. With markings, a contracted leaf might have one node and two markings, and would then be stable. Thus the absence of markings is an essential hypothesis in the argument.


\section{The cuspidal roadmap}
\MCsolution{ex:cusp-roadmap}
The isomorphism-sheaf argument glues underlying morphisms,
their inverses, and the compatibility equations. Effective
object descent then uses the functorial ample bundle to
descend the total spaces, followed by the markings and the
map, and the assumed locality to recover all defining
conditions. Along with the category and base-change checks,
this proves the stack axioms.

It does not construct an algebraic atlas, prove boundedness,
represent the diagonal, exclude infinitesimal automorphisms,
or prove existence and uniqueness of limits. Those are
needed for the proposed proper Deligne--Mumford conclusion.
The marked cuspidal cubic has an ample $\omega_C(p)$ but a
positive-dimensional automorphism group, demonstrating that
even the automorphism part does not follow from ampleness
in the cuspidal setting.
